\subsection{Invariant backgrounds on compact Lie-group quotients}
\label{sec:parallelizable-gauduchon-cones}

A normalized invariant volume reduces the Gauduchon quotient to a matrix determinant. The resulting metric controls the eigenvalues of the normalized minimizing form. This gives complete cones of nonclosed backgrounds, with an explicit distinction between the coefficient boundary and the unit-volume completion.

Let $G$ be a connected complex Lie group of dimension $n\geq2$, and let $X=\Gamma\backslash G$ be compact, where the discrete subgroup $\Gamma$ acts on the left. Choose a left-invariant holomorphic coframe $\alpha_1,\ldots,\alpha_n$, which descends to $X$. For a Hermitian matrix $A$ put
\[
 \begin{gathered}
 \theta_A=\frac{\mathrm i}{2\pi}\sum_{j,k}A_{jk}\alpha_j\wedge\bar\alpha_k,
 \qquad h_0=\theta_I,\\
 H=\int_Xh_0^n,\qquad \mu=\frac{h_0^n}{H}.
 \end{gathered}
\]
Fix a real linear subspace $W\subset\operatorname{Herm}(n)$ containing a positive matrix, and write $V_W=\{\theta_A:A\in W\}$. Here the probability volume is the normalized invariant volume just specified.

\Needspace{12\baselineskip}
\begin{proposition}\label{prop:parallelizable-gaud-volume}\index{Gauduchon variational volume!invariant backgrounds}
The backgrounds in $V_W$ admitting a smooth positive representative after a smooth $\ddc$ change are precisely $\theta_A$ with $A>0$. Moreover,
\[
 V_W\cap\ddc C^\infty(X,\R)=\{0\}.
\]
Every positive $\theta_A$ is balanced, and its unique minimizing Gauduchon metric normalized by $D_\mu=1$ is
\begin{equation}\label{eq:parallelizable-gaud-volume}
 \lambda_\mu(\theta_A)=(H\det A)^{1/n},\qquad
 \gamma_A=\frac{\theta_A}{\lambda_\mu(\theta_A)},\qquad
 \gamma_A^n=\mu.
\end{equation}
On this coefficient cone the Hessian metric is
\begin{equation}\label{eq:parallelizable-gaud-hessian}
 g_A(U,Z)=\frac1n\operatorname{tr}(A^{-1}UA^{-1}Z),
 \qquad U,Z\in W.
\end{equation}

For $A\in W$, a positive-current representative $\theta_A+\ddc u$, with a real distributional potential $u$, exists exactly when $A\geq0$. The coefficient closure is $W\cap\{A\geq0\}$, and its canonical representative is the smooth semipositive form $\theta_A$. The zero background has only the zero positive-current representative. If $G$ is nonabelian, every positive background in this cone is nonclosed and has no closed representative obtained by adding a smooth $\ddc$ potential.
\end{proposition}

\begin{proof}
\emph{Invariant averages and balancedness.}
Let $Z_j$ be dual to $\alpha_j$, and set $\Upsilon=\alpha_1\wedge\cdots\wedge\alpha_n$. This holomorphic top form and its conjugate are closed. For an invariant holomorphic vector field $Z$, evaluation on the invariant frame gives
\[
 \mathcal L_Z\Upsilon=-\operatorname{tr}(\operatorname{ad}Z)\Upsilon,
 \qquad \mathcal L_Z\bar\Upsilon=0.
\]
The second identity also follows from Cartan's formula. Since $h_0^n$ is a constant multiple of $\Upsilon\wedge\bar\Upsilon$, with the positive orientation factor, compact Stokes gives
\[
 0=\int_X\mathcal L_Zh_0^n=-H\operatorname{tr}(\operatorname{ad}Z).
\]
Thus the fields and their conjugates have zero divergence. In particular,
\begin{equation}\label{eq:parallelizable-zero-mean}
 \int_XZ_jf\,\mu=\int_X\bar Z_jf\,\mu=0
 \qquad(f\in C^\infty(X,\mathbb C)).
\end{equation}
Cartan's formula also gives $\mathrm d(\iota_{Z_j}\Upsilon)=0$. These contractions span the invariant $(n-1,0)$-forms. The form $\theta_A^{n-1}$ is a constant linear combination of their wedges with conjugates, so $\mathrm d\theta_A^{n-1}=0$.

The holomorphic fields commute with the antiholomorphic fields, and $\mathrm d\alpha_j$ has type $(2,0)$. The coefficient matrix of $\ddc u$ is therefore $(Z_j\bar Z_k u)$. For smooth $u$ its entries have mean zero by \eqref{eq:parallelizable-zero-mean}. The mean matrix of $\theta_A+\ddc u$ is $A$. Smooth positivity implies positivity of the mean quadratic form on every nonzero constant vector, hence $A>0$. The converse follows by taking $u=0$. The same average proves the absence of nonzero smooth $\ddc$ directions in $V_W$.

For distributions the adjoints with respect to $\mu$ annihilate constants, so the mean of each $Z_j\bar Z_k u$ is still zero. The coefficients of a positive current in this global frame are a positive matrix of measures. Their mean is semipositive, giving $A\geq0$. Conversely $\theta_A$ itself is positive when $A\geq0$. If $A=0$, the positive trace measure has total mass zero; polarization then makes the entire current zero. Finally, for a fixed positive $A_+\in W$, the matrices $A+\varepsilon A_+$ approximate every semipositive $A\in W$ from the positive cone. This proves the closure assertion.

\emph{The quotient and its equality case.}
Put $Q=H\det A$. The determinant identity is $\theta_A^n=Q\mu$. For every smooth positive Hermitian form $\gamma$, the matrix arithmetic--geometric mean inequality gives
\[
 \frac{\theta_A\wedge\gamma^{n-1}}{\gamma^n}
 =\frac1n\operatorname{tr}_\gamma\theta_A
 \geq\left(\frac{\theta_A^n}{\gamma^n}\right)^{1/n}.
\]
Integration yields $\int_X\theta_A\wedge\gamma^{n-1}\geq Q^{1/n}D_\mu(\gamma)$. This bounds every Gauduchon test in \cref{def:gaud-variational-volume}. The balanced test $\gamma=\theta_A$ has $D_\mu(\theta_A)=Q^{(n-1)/n}$ and attains equality, proving the root-volume formula.

At a minimizing smooth Gauduchon test the continuous nonnegative deficit has integral zero. Equality pointwise forces $\gamma=f\theta_A$ for a smooth positive $f$. Set $q=f^{n-1}$ and $T=\theta_A^{n-1}$. Since $\mathrm dT=0$, the Gauduchon equation becomes $\ddc q\wedge T=0$, and
\[
 0=\int_Xq\ddc q\wedge T=-\int_X\mathrm dq\wedge\dc q\wedge T.
\]
The last integrand is nonnegative and vanishes exactly where $\mathrm dq=0$. Thus $f$ is constant. As $n\geq2$ and
\[
 D_\mu(c\theta_A)=c^{n-1}\lambda_\mu(\theta_A)^{n-1},
\]
normalization determines $c=\lambda_\mu(\theta_A)^{-1}$. This proves the unique minimizer and its volume identity.

Differentiating $\log\lambda_\mu(\theta_A)=(\log H+\log\det A)/n$ gives \eqref{eq:parallelizable-gaud-hessian}. Its square norm is $\|A^{-1/2}UA^{-1/2}\|_{\mathrm F}^2/n$, which is positive for $U\ne0$. Here $\|B\|_{\mathrm F}^2=\operatorname{tr}(B^*B)$ is the Frobenius norm.

Finally,
\[
 \partial\theta_A=\frac{\mathrm i}{2\pi}\sum_{j,k}A_{jk}\,\mathrm d\alpha_j\wedge\bar\alpha_k.
\]
For invertible $A$, vanishing forces every $\mathrm d\alpha_j=0$, by independence of the conjugate coframe. The Lie algebra would then be abelian. Addition of $\ddc u$ does not change $\mathrm d\theta_A$, proving the last assertion.
\end{proof}

\Needspace{12\baselineskip}
\begin{lemma}\label{lem:positive-matrix-hessian-distance}\index{Gauduchon variational volume!matrix distance}
The full positive Hermitian matrix cone, with the metric in \eqref{eq:parallelizable-gaud-hessian}, has distance
\begin{equation}\label{eq:positive-matrix-hessian-distance}
 d_{\mathrm{full}}(A,B)=\frac1{\sqrt n}
       \|\log(A^{-1/2}BA^{-1/2})\|_{\mathrm F}.
\end{equation}
For endpoints with equal determinant, this is also the intrinsic distance in their full fixed-determinant hypersurface.
\end{lemma}

\begin{proof}
Fixed congruence is an isometry by the trace formula. Reduce to endpoints $I,C>0$. The curve $\exp(t\log C)$ has constant speed $\|\log C\|_{\mathrm F}/\sqrt n$.

For the reverse inequality, diagonalization gives the convergent resolvent formula and its derivative
\begin{align*}
 \log M&=\int_0^\infty\bigl((1+s)^{-1}I-(M+sI)^{-1}\bigr)\,\mathrm ds,\\
 D\log_M[U]&=\int_0^\infty(M+sI)^{-1}U(M+sI)^{-1}\,\mathrm ds.
\end{align*}
On a neighborhood with a positive spectral lower bound, differentiation of the inverse is dominated by an integrable function. In a unitary basis diagonalizing $M$, with eigenvalues $d_i>0$, the derivative has entries $c_{ij}U_{ij}$, where
\[
 c_{ij}=\frac{\log d_i-\log d_j}{d_i-d_j}\quad(d_i\ne d_j),
 \qquad c_{ij}=\frac1{d_i}\quad(d_i=d_j).
\]
Writing $d_i/d_j=\mathrm e^{2s}$ gives $\sqrt{d_id_j}|c_{ij}|=|s/\sinh s|\leq1$, including the limiting value at $s=0$. Hence
\[
 \|D\log_M[U]\|_{\mathrm F}
 \leq\|M^{-1/2}UM^{-1/2}\|_{\mathrm F}.
\]
For any piecewise $C^1$ positive matrix path from $I$ to $C$, integration bounds its length below by $\|\log C\|_{\mathrm F}/\sqrt n$. This proves the distance formula. Restoring the endpoints, the attaining curve is
\begin{equation}\label{eq:positive-matrix-exponential-path}
 M(t)=A^{1/2}(A^{-1/2}BA^{-1/2})^tA^{1/2}.
\end{equation}
Its determinant is $(\det A)^{1-t}(\det B)^t$. Thus for equal-determinant endpoints it stays in the indicated hypersurface.
\end{proof}

Write $\lambda(A)=\lambda_\mu(\theta_A)$ and
\[
 S_W=\{A\in W:A>0,\ H\det A=1\}.
\]
The radial derivative of $\lambda$ is positive, so $S_W$ is a smooth hypersurface of the coefficient cone. All distances $d_W$ below are intrinsic path distances on this fixed hypersurface with its induced Hessian metric.

\Needspace{12\baselineskip}
\begin{theorem}\label{thm:parallelizable-unit-volume-completion}\index{Gauduchon variational volume!complete invariant cone}
The space $(S_W,d_W)$ is complete and proper. A sequence is Cauchy exactly when its coefficients converge to a positive matrix in $S_W$. For $A,B\in S_W$,
\begin{align}
 d_W(A,B)&\geq\frac1{\sqrt n}
       \|\log(A^{-1/2}BA^{-1/2})\|_{\mathrm F},\label{eq:parallelizable-intrinsic-lower}\\
 d_W(A,B)&\leq\frac{\|A-B\|_{\mathrm F}}
       {\sqrt n\min\{\lambda_{\min}(A),\lambda_{\min}(B)\}}.
       \label{eq:parallelizable-intrinsic-upper}
\end{align}
The actual normalized minimizing forms satisfy pointwise on $X$
\begin{equation}\label{eq:parallelizable-minimizer-order}
 \mathrm e^{-\sqrt n d_W(A,B)}\gamma_A
 \leq\gamma_B\leq\mathrm e^{\sqrt n d_W(A,B)}\gamma_A.
\end{equation}
Thus every completion point has its canonical smooth positive current representative, and the completion adds no semipositive coefficient boundary point.

If $A_j>0$ in $W$ converges to a nonzero singular semipositive $A_0$, then $\widehat A_j=A_j/\lambda(A_j)$ lies in $S_W$ and, for fixed $A_\star\in S_W$,
\begin{equation}\label{eq:parallelizable-singular-boundary-rate}
 d_W(A_\star,\widehat A_j)
 \geq\frac{|\log\det A_j|}{n\sqrt n}+O(1)
 \longrightarrow\infty.
\end{equation}
Normalization of a sequence tending to zero can instead converge to an existing positive point of $S_W$.
\end{theorem}

\begin{proof}
\emph{Paths in the specified subspace.}
Every path in $S_W$ is a path in the full fixed-determinant hypersurface, so \cref{lem:positive-matrix-hessian-distance} gives \eqref{eq:parallelizable-intrinsic-lower}. To obtain an upper bound within $W$, put
\[
 U=B-A,\qquad C(t)=A+tU,\qquad N(t)=\frac{C(t)}{\lambda(C(t))}.
\]
The positive cone is convex and $W$ is linear, so $N(t)$ stays in $S_W$ and joins $A$ to $B$. If $K=C^{-1/2}UC^{-1/2}$, differentiation gives
\[
 \begin{gathered}
 g_N(N',N')=\frac1n\operatorname{tr}(K^2)
                 -\frac1{n^2}(\operatorname{tr}K)^2
 \leq\frac{\|U\|_{\mathrm F}^2}{nm^2},\\
 m=\min\{\lambda_{\min}(A),\lambda_{\min}(B)\}.
 \end{gathered}
\]
Indeed $C(t)\geq mI$ and $(\log\lambda(C))'=\operatorname{tr}(C^{-1}U)/n$. Integrating proves \eqref{eq:parallelizable-intrinsic-upper}, and also connectedness of $S_W$.

Put $d=d_W(A,B)$. The logarithmic spectral lower bound gives
\begin{equation}\label{eq:parallelizable-coefficient-order}
 \mathrm e^{-\sqrt n d}A\leq B\leq\mathrm e^{\sqrt n d}A.
\end{equation}
Since $\gamma_A=\theta_A$ on $S_W$, this proves the pointwise minimizing-form estimate.

\emph{Cauchy sequences and bounded balls.}
For an intrinsic Cauchy sequence, comparison with one fixed tail element gives $mI\leq A_j\leq MI$ with fixed positive $m,M$. Comparing two tail elements then gives
\[
 \|A_j-A_k\|_{\mathrm{op}}
 \leq M\bigl(\mathrm e^{\sqrt n d_W(A_j,A_k)}-1\bigr).
\]
Its coefficients are Cauchy. Their limit lies in the closed linear space $W$, remains positive, and has $H\det A=1$. The intrinsic upper bound proves convergence to that point. Conversely coefficient convergence to a positive point gives intrinsic convergence by the same bound. This proves completeness and the exact Cauchy criterion. The minimizing forms converge in $C^\infty$, because their coefficients multiply a fixed smooth coframe.

The two topologies agree by these estimates. A closed ball of radius $R$ around $A_\star$ is contained in the compact matrix interval
\[
 \begin{gathered}
 \mathrm e^{-\sqrt nR}A_\star\leq A\leq\mathrm e^{\sqrt nR}A_\star,\\
 A\in W,\qquad H\det A=1.
 \end{gathered}
\]
It is closed in that interval, hence compact. This proves properness. Distances from a fixed reference are bounded exactly when the relative eigenvalues are bounded above and away from zero.

\emph{The coefficient boundary.}
For $A_j\to A_0\geq0$ singular and nonzero, $\lambda(A_j)\to0$ while $\lambda_{\max}(A_j)$ has a positive limit. The largest eigenvalue of $\widehat A_j$ tends to infinity. Applying \eqref{eq:parallelizable-coefficient-order} to $A_\star,\widehat A_j$ gives
\[
 \begin{aligned}
 d_W(A_\star,\widehat A_j)
 &\geq\frac1{\sqrt n}\log
       \frac{\lambda_{\max}(A_j)}{\lambda(A_j)\lambda_{\max}(A_\star)}\\
 &=\frac{|\log\det A_j|}{n\sqrt n}+O(1).
 \end{aligned}
\]
The error is bounded because $A_0\ne0$. This proves \eqref{eq:parallelizable-singular-boundary-rate}. For comparison, $A_j=t_jA_\star$ with $t_j\downarrow0$ normalizes identically to $A_\star$. Its raw zero-current limit supplies no new point of the unit-volume completion.
\end{proof}

By \cref{prop:gaud-volume-metric-product}, the full prescribed cone is complete as well. The invariant probability and the fixed linear subspace are the hypotheses used here; other geometries can have the finite-distance boundary of \cref{ex:gaud-metric-blowup}.

The exponential path need not stay in a proper $W$. For $n=3$, take
\[
 W=\operatorname{span}_{\R}\{I,\operatorname{diag}(0,1,2)\},
 \qquad a=H^{-1/3}.
\]
Its unit-volume endpoints $aI$ and $a6^{-1/3}\operatorname{diag}(1,2,3)$ have full-matrix midpoint $a6^{-1/6}\operatorname{diag}(1,\sqrt2,\sqrt3)$. The diagonal entries fail the relation $x_1-2x_2+x_3=0$, since $1-2\sqrt2+\sqrt3\ne0$. Thus this path leaves $W$, whereas the normalized chord in the proof remains in it.

\Needspace{12\baselineskip}
\begin{example}\label{ex:iwasawa-complete-gaud-cone}\index{Iwasawa manifold!Gauduchon background cone}
On $G=\mathbb C^3$ take the multiplication
\[
 z\cdot w=(z_1+w_1,z_2+w_2,z_3+w_3+z_1w_2),
\]
the Gaussian lattice $\Gamma=\mathbb Z[\mathrm i]^3$, and $X=\Gamma\backslash G$. This is a compact non-K\"ahler manifold of dimension three, with invariant coframe
\[
 \alpha_1=\mathrm dz_1,\qquad \alpha_2=\mathrm dz_2,\qquad
 \alpha_3=\mathrm dz_3-z_1\mathrm dz_2.
\]
For its full nine-dimensional cone $W=\operatorname{Herm}(3)$ and the normalized invariant probability,
\[
 H=\frac6{\pi^3},\qquad
 \mathcal V_\mu(\theta_A)=\frac6{\pi^3}\det A,\qquad
 S=\{A>0:\det A=\pi^3/6\}.
\]
Every positive background is nonclosed. The eight-dimensional unit-volume hypersurface is complete and proper, with exact intrinsic distance
\begin{equation}\label{eq:iwasawa-unit-volume-distance}
 d_S(A,B)=\frac1{\sqrt3}\|\log(A^{-1/2}BA^{-1/2})\|_{\mathrm F}.
\end{equation}
\end{example}

\begin{proof}
The bilinear cross term makes the multiplication associative, with identity zero and inverse $(-z_1,-z_2,-z_3+z_1z_2)$. These operations preserve the Gaussian lattice. Its action is free and properly discontinuous: elements with $\gamma K\cap K\ne\varnothing$ belong to the compact set $KK^{-1}$, containing only finitely many lattice elements.

Choose Gaussian integers $\gamma_1,\gamma_2$ to reduce $z_1,z_2$ to their half-open unit squares, then $\gamma_3$ to reduce $z_3+\gamma_1z_2$. This gives a representative in the real six-dimensional unit cube. Two representatives in the half-open cube force $\gamma_1=\gamma_2=0$, then $\gamma_3=0$. The cube is therefore a measurable fundamental domain, and the image of its closure proves compactness.

The displayed coframe is left-invariant, descends to $X$, and has
\[
 \mathrm d\alpha_3=-\alpha_1\wedge\alpha_2,\qquad
 \alpha_1\wedge\alpha_2\wedge\alpha_3
       =\mathrm dz_1\wedge\mathrm dz_2\wedge\mathrm dz_3.
\]
The density of $h_0^3$ is $6/\pi^3$ times real Lebesgue volume on the fundamental cube, giving the stated $H$. The coefficient of $\alpha_1\wedge\alpha_2\wedge\bar\alpha_3$ in $\partial\theta_A$ is $-\mathrm iA_{33}/(2\pi)$, so positivity makes every such background nonclosed.

To exclude an arbitrary K\"ahler metric, suppose a positive closed $\kappa$ existed and average its coefficient functions in this coframe against $\mu$. Their mean is a positive matrix $A$. In the coefficients of $\partial\kappa$, derivative terms have mean zero by \eqref{eq:parallelizable-zero-mean}; the remaining means are exactly the coefficients of $\partial\theta_A$. Closedness would force them to vanish, contradicting $A_{33}>0$. Thus $X$ is non-K\"ahler.

All volume, minimizer and completeness assertions follow from \cref{prop:parallelizable-gaud-volume,thm:parallelizable-unit-volume-completion}. As $W$ is the full matrix space, the equal-determinant exponential path stays in $S$, and \cref{lem:positive-matrix-hessian-distance} gives \eqref{eq:iwasawa-unit-volume-distance}. For $A_\star=aI$, $a=(\pi^3/6)^{1/3}$, this becomes
\[
 d_S(A_\star,A)^2=\frac13\sum_{j=1}^3
       \left(\log\frac{\lambda_j(A)}a\right)^2.
\]
The three logarithms sum to zero.

For a concrete degeneration take $A_\varepsilon=\operatorname{diag}(1,1,\varepsilon)$, $0<\varepsilon\leq1$. Its coefficient limit gives the nonzero smooth semipositive current $\theta_{\operatorname{diag}(1,1,0)}$. The normalized minimizing matrices and their distances are instead
\begin{equation}\label{eq:iwasawa-diagonal-degeneration}
 \begin{gathered}
 \widehat A_\varepsilon
   =a\operatorname{diag}(\varepsilon^{-1/3},\varepsilon^{-1/3},\varepsilon^{2/3}),\\
 d_S(A_\star,\widehat A_\varepsilon)=\frac{\sqrt2}{3}|\log\varepsilon|.
 \end{gathered}
\end{equation}
The coefficient boundary current therefore lies at infinite normalized parameter distance.
\end{proof}
